\documentclass[11pt,a4paper]{article}
\usepackage[margin=2.2cm]{geometry}
\usepackage{booktabs,array,hyperref,enumitem}
\usepackage[T1]{fontenc}
\usepackage{lmodern}
\setlist{itemsep=2pt,topsep=3pt}
\hypersetup{colorlinks=true,urlcolor=blue}
\setlength{\parindent}{0pt}
\setlength{\parskip}{5pt}

\title{RAG Assignment: Baseline Retrieval vs.\ Cross-Encoder Re-Ranking}
\author{Pavan Reddy}
\date{}

\begin{document}
\maketitle

\textbf{Code:} \url{https://github.com/pavanreddy2480/RAG_LMA} (single notebook \texttt{rag\_assignment.ipynb}, with outputs saved).

\section*{Setup}
The document is the FIDE Laws of Chess (25 pages). It has many exceptions, penalties and ``not allowed'' clauses, which suits the stress tests. I split it into 121 chunks of about 100 words with a 20-word overlap; both pipelines use these same chunks.
\begin{itemize}
  \item Bi-encoder: \texttt{all-MiniLM-L6-v2}, normalised vectors, FAISS inner-product index (equal to cosine).
  \item Cross-encoder: \texttt{cross-encoder/ms-marco-MiniLM-L-6-v2}.
  \item LLM: \texttt{gemini-3.1-flash-lite-preview}, temperature 0, one shared prompt for both answers.
  \item Pipeline 1: top-10 from the bi-encoder, top-3 to the LLM. Pipeline 2: the same 10 candidates re-sorted by the cross-encoder, top-3 to the LLM.
\end{itemize}

\section{Architecture Comparison}
A bi-encoder encodes the query and each chunk separately into one vector each, and the score is the cosine between the two vectors. Because the chunk vectors do not depend on the query, they are computed once and stored in an index. The price is that a whole chunk is squeezed into a single vector before the query is seen.

A cross-encoder takes the query and the chunk together as one input, so self-attention runs across all query tokens and chunk tokens jointly, and it outputs one relevance score. It can therefore match specific words in the query to specific words in the chunk (for example ``always'' in the query against ``cannot win\ldots draw'' in the chunk). Nothing can be precomputed, because the score depends on the query.

\section{Qualitative and Rank Shift Analysis}
For each query I marked the ``gold'' chunk by a phrase from the clause that really answers it, found by reading the rulebook before running the models. Table~\ref{tab:ranks} gives the rank of that chunk among the 10 candidates before and after re-ranking.

\textbf{How the queries were chosen.} I wrote more than 40 candidate queries and checked, using retrieval only, where the gold chunk landed. The queries below are the ones that showed the behaviour I wanted to test, plus a few controls. So this is a set of examples chosen to expose failures, not a fair estimate of how often re-ranking helps. Many negation and exception queries I tried were already answered correctly by the baseline (gold chunk at rank 1), and these are not listed.

\begin{table}[h]
\centering\small
\begin{tabular}{@{}p{6.4cm}p{3.1cm}ccc@{}}
\toprule
Query (type) & Gold chunk & Bi & CE & Top-3 before $\to$ after \\
\midrule
Extra minutes when the arbiter rejects a draw claim (buried chunk) & 63 & 4 & 1 & No $\to$ Yes \\
If my opponent's time runs out, do I always win? (exception) & 66 (and 65) & 10 & 1 & No $\to$ Yes \\
Does an illegal move always cost the player? (exception) & 45 & 5 & 3 & No $\to$ Yes \\
Can I still claim a draw after touching a piece? (negation) & 58 & 9 & 1 & No $\to$ Yes \\
Can the scoresheet be used for personal notes? (negation) & 66 & 6 & 1 & No $\to$ Yes \\
What can the referee do to a player who breaks the rules? (synonym) & 69 & 9 & 2 & No $\to$ Yes \\
How long can a player be absent before the game is forfeited? (synonym) & 33 (and 32) & 5 & 9 & No $\to$ No \\
Third illegal move (control) & 45 & 1 & 1 & Yes $\to$ Yes \\
Withdraw a draw claim (negation control) & 59 & 1 & 1 & Yes $\to$ Yes \\
\bottomrule
\end{tabular}
\caption{Rank of the gold chunk before (bi-encoder) and after (cross-encoder) re-ranking.}
\label{tab:ranks}
\end{table}

\textbf{Exceptions and corner cases.} For ``do I always win?'' the clause with the exception (``if the opponent cannot win the game by any series of legal moves, his score shall be a draw'') sat at rank 10, so the baseline never gave it to the LLM. Answer 1 was \emph{``Not found in the provided context''}; the cross-encoder moved the chunk to rank 1 and Answer 2 stated the draw exception correctly. For ``does an illegal move always cost the player?'' the baseline top-3 (chunks 25, 22, 78) did not contain the clause on the third illegal move (chunk 45, rank 5), and Answer 1 was again ``Not found''. After re-ranking chunk 45 was third, and Answer 2 correctly said that the first two illegal moves only give the opponent two extra minutes, the game is lost on the third, and a draw applies if the opponent cannot checkmate.

\textbf{Negation and constraints.} ``Can I still claim a draw after touching a piece?'' needs the clause saying the player \emph{loses} the right to claim (chunk 58). The bi-encoder ranked it 9th, behind chunks about claiming draws in general, and Answer 1 was ``Not found''. The cross-encoder put it first, and Answer 2 said that touching a piece without having claimed the draw loses the right to claim on that move. For the scoresheet query the gold chunk also moved from 6 to 1, but both answers were ``Not found'', so the LLM did not extract the rule even with the right chunk present. This shows that a better ranking does not guarantee a better answer. Several other negation queries (for example ``not allowed to withdraw a claim'', used as a control above) were already ranked first by the bi-encoder, so the bi-encoder does not always blur negation.

\textbf{Keyword trap / synonyms.} For ``referee'' the rulebook says ``arbiter'', so the penalty list (chunk 69) was ranked 9th by the bi-encoder and 2nd by the cross-encoder. Both answers listed the penalties, but Answer 2 also added the penalty for persistent refusal to comply, which Answer 1 left out. The opposite happened for ``absent\ldots forfeited'': the default-time clause (chunk 33) fell from rank 5 to 9 and both answers described the adjourned-game rule (``more than one hour late'') instead. I think the MS-MARCO cross-encoder preferred chunks containing words like ``late'' and ``loses'', but I did not test this further.

\textbf{Buried chunk.} For ``extra minutes'' the chunk with the exact sentence moved from rank 4 to 1 (baseline top-3 was 62, 61, 59; re-ranked 63, 62, 61). Both answers were correct (``two extra minutes''), because chunk 62 contains a similar clause about postponed decisions. So the ranking improved, but the baseline did not fail.

\textbf{Summary.} The gold chunk entered the top-3 for 6 of 9 queries. In 3 of these (``do I always win'', ``illegal move'', ``claim a draw after touching'') the baseline answered ``Not found'' and the re-ranked pipeline answered correctly. In the synonym query the re-ranked answer was slightly more complete. In one query (``absent'') re-ranking made retrieval worse, and in the controls nothing changed. No baseline answer contained an invented rule: with the instruction to answer only from the context, the baseline failed by omission (``Not found'') rather than by hallucination. Because the queries were selected as described above, these counts are not a benchmark.

\section{Latency Profiling}
Retrieval timings are means over the 9 queries with 10 repeats each, on a laptop, with 10 candidates. The LLM time is one uncached API call, since it is identical for both pipelines and the free API quota is small.

\begin{table}[h]
\centering\small
\begin{tabular}{@{}lr@{}}
\toprule
Stage & Time \\
\midrule
Pipeline 1: bi-encoder retrieval (embed query + FAISS top-10) & 8.1 ms \\
Re-ranking step (10 cross-encoder pairs) & 28.2 ms \\
Pipeline 2 retrieval total (bi-encoder + cross-encoder) & 36.3 ms (4.5$\times$ baseline) \\
LLM generation (same in both pipelines) & about 0.7 s \\
\midrule
Cross-encoder on all 121 chunks, one query (for comparison) & 381 ms \\
\bottomrule
\end{tabular}
\caption{Runtime of the two pipelines.}
\end{table}

Re-ranking makes retrieval about 4.5 times slower, but the absolute cost is about 28 ms, which is small next to the LLM call.

\section{Scaling Question}
A cross-encoder must run a full transformer pass for every (query, chunk) pair, and the result cannot be cached because it depends on the query. Running it over the whole collection costs a number of passes proportional to the number of chunks for every query. Here 121 chunks already cost 381 ms, about 14 times the cost of scoring 10 chunks, and a collection of millions of chunks would take minutes per query.

A bi-encoder avoids this: chunk vectors are computed once offline, and a query costs one encoding plus a fast nearest-neighbour search. So the first stage uses the cheap model to get a short list with good recall, and the expensive cross-encoder only refines that short list. The limit of this design is that if the right chunk is not in the top-10, the re-ranker cannot recover it.

\newpage
\appendix

\section{Full Results per Query}

For every query: the complete re-ordering of the 10 candidates (chunk id, bi-encoder rank $\to$ cross-encoder rank; the chunk with $^*$ holds the gold clause) and both LLM answers, copied from the notebook output.\par

\subsection*{Q1. How many extra minutes does the opponent get when the arbiter rejects a draw claim?}
\emph{Failure query 1 (gold chunk \#4 -> \#1)}\par
{\small chunk: bi$\to$CE rank\par\smallskip \begin{tabular}{@{}lllll@{}}
63$^*$: 4$\to$1 & 62: 1$\to$2 & 61: 2$\to$3 & 59: 3$\to$4 & 113: 8$\to$5 \\
66: 9$\to$6 & 53: 6$\to$7 & 36: 7$\to$8 & 120: 10$\to$9 & 52: 5$\to$10
\end{tabular}}\par\smallskip
\textbf{Answer 1 (baseline):} The opponent shall be awarded two extra minutes.\par
\textbf{Answer 2 (re-ranked):} The opponent shall be awarded two extra minutes.\par

\subsection*{Q2. What happens if a player makes a third illegal move?}
\emph{Control (easy query, no shift expected)}\par
{\small chunk: bi$\to$CE rank\par\smallskip \begin{tabular}{@{}lllll@{}}
45$^*$: 1$\to$1 & 22: 6$\to$2 & 75: 4$\to$3 & 25: 2$\to$4 & 58: 10$\to$5 \\
44: 3$\to$6 & 78: 5$\to$7 & 26: 7$\to$8 & 27: 8$\to$9 & 52: 9$\to$10
\end{tabular}}\par\smallskip
\textbf{Answer 1 (baseline):} For a third illegal move by the same player, the arbiter shall declare the game lost by that player. However, the game is drawn if the position is such that the opponent cannot checkmate the player’s king by any possible series of legal moves.\par
\textbf{Answer 2 (re-ranked):} For a third illegal move by the same player, the arbiter shall declare the game lost by that player. However, the game is drawn if the position is such that the opponent cannot checkmate the player’s king by any possible series of legal moves.\par

\subsection*{Q3. How long can a player be absent before the game is forfeited?}
\emph{Keyword trap / synonyms (re-ranker does not help)}\par
{\small chunk: bi$\to$CE rank\par\smallskip \begin{tabular}{@{}lllll@{}}
117: 1$\to$1 & 63: 2$\to$2 & 118: 3$\to$3 & 114: 10$\to$4 & 62: 8$\to$5 \\
116: 6$\to$6 & 110: 7$\to$7 & 119: 9$\to$8 & 33$^*$: 5$\to$9 & 52: 4$\to$10
\end{tabular}}\par\smallskip
\textbf{Answer 1 (baseline):} A player loses the game if they arrive at the chessboard more than one hour late for the resumption of an adjourned game, unless the competition rules or the arbiter specify otherwise.\par
\textbf{Answer 2 (re-ranked):} A player loses the game if they arrive at the chessboard more than one hour late for the resumption of an adjourned game, unless the competition rules or the arbiter specify otherwise.\par

\subsection*{Q4. If my opponent's time runs out, do I always win?}
\emph{Failure query 2 + Exception / corner case}\par
{\small chunk: bi$\to$CE rank\par\smallskip \begin{tabular}{@{}lllll@{}}
66$^*$: 10$\to$1 & 63: 2$\to$2 & 33: 1$\to$3 & 61: 3$\to$4 & 34: 6$\to$5 \\
59: 4$\to$6 & 37: 8$\to$7 & 110: 9$\to$8 & 29: 5$\to$9 & 32: 7$\to$10
\end{tabular}}\par\smallskip
\textbf{Answer 1 (baseline):} Not found in the provided context.\par
\textbf{Answer 2 (re-ranked):} If the opponent cannot win the game by any series of legal moves, the score shall be a draw.\par


\end{document}
